\section{Preliminaries}

\subsection{Task Definition} \label{sec:task_definition}

We investigate the problem of reinforcement learning in LLM post-training stage for complex, mixed-domain text generation tasks. The core objective is to train a single policy model that achieves high performance across multiple task domains without sacrificing the quality of any individual task.
Let \(\{ D_1, D_2, \dots, D_n \}\) denote datasets from \(n\) different domains, each corresponding to one type of task (e.g., translation, summarization, question answering).  
We define a mixed-domain distribution \(D = \{ D_1, D_2, \dots, D_n \}\).
Our goal is to train a policy model \(A\) that maximizes expected performance over \(D\) under a multi-task reinforcement learning framework:  
\[
\max_{A} \; \mathbb{E}_{d \sim D} \left[ R(A(d)) \right],
\]  
where \(R(A(d))\) denotes the reward of policy inference on sample \(d\). The reward is computed from test-time metrics, such as success rate, accuracy, LLM-judge score, and other task-specific metrics. 
We study the blended domains where real LLM post-training usually focuses: \textbf{mathematical reasoning}, \textbf{code and programming}, \textbf{translation}, and \textbf{general conversation} tasks. We show information about the datasets used in our experiments in Table~\ref{tab:dataset_stat}.


\subsection{LLM post-training with online RL} \label{sec:rl_algo}


In reinforcement learning from human feedback (RLHF), the \textbf{Proximal Policy Optimization} algorithm~\citep{ppo} is frequently employed for policy optimization.  
The typical workflow begins with a \textit{warm-start training} phase in which a \textbf{Value Model} (often a reward model) is learned. Training triplets \((x, y_{+}, y_{-})\) are sampled from human preference data, with \(y_{+}\) being preferred over \(y_-\). Let \(r_{\theta}(x, y)\) be the reward assigned by the model given prompt \(x\) and response \(y\); the learning objective function can be expressed as:
\begin{equation*}
    \mathcal{L}(\theta) = -\frac{1}{N} \,\mathbb{E}_{(x,y_+, y_-) \sim D } \bigg[ \log \sigma \big( \Delta r_{\theta} \big) \bigg],
\end{equation*}
where \(\sigma(\cdot)\) is the logistic sigmoid function and \(\Delta r_{\theta} = r_{\theta}(x, y_+) - r_{\theta}(x, y_-)\)  .

We adopt a more training-efficient framework. The \textbf{Group Relative Policy Optimization} (GRPO)~\citep{grpo} modifies the advantage estimation to reduce the dependence on a learnable value model for estimating the advantage baseline. Instead, GRPO computes the normalized advantage within a group of sampled outputs \(\hat{A}_{i} = \frac{r_i - \mathrm{mean}(r)}{\mathrm{std}(r)}  \), 
where \(r=\{ r_i \}_{i=1}^G\) are the rewards assigned to \(G\) candidate outputs for the same prompt, \(\mathrm{mean}(r)\) and \(\mathrm{std}(r)\) are computed over the group. Also, the Kullback-Leibler Divergence term is removed from the per-step reward and is instead applied directly to the overall optimization objective:
\begin{align}
\max_{\phi} \; \;
& \mathbb{E}_{\substack{x \sim D, \\ \{y_i\}_{i=1}^G \sim \pi_{\mathrm{old}}(y_i|x)}} 
\Bigg[
   \frac{1}{G} \sum_{i=1}^G
   \min \Bigg\{
      \frac{\pi_{\theta}(y_i|x)}{\pi_{\mathrm{old}}(y_i|x)} ,  \;
      \nonumber \\[-0.5em]
& 
      \mathrm{clip}\!\left(
         \frac{\pi_{\theta}(y_i|x)}{\pi_{\mathrm{old}}(y_i|x)},
         1 - \epsilon, \;
         1 + \epsilon
      \right) 
   \Bigg\} \hat{A}_{i},
\Bigg] 
- \beta \, \mathbb{D}_{\mathrm{KL}}\!\left[ \pi_{\phi} \;\|\; \pi_{\mathrm{ref}} \right].
\end{align}

% % 插入一个表格，让读者一目了然
\begin{table}[t]
    \centering
    \caption{Summary of domains considered across Math, Code, Translation, and General Conversation.}
    \label{tab:dataset_stat}
    \small % Slightly smaller font to ensure fit if necessary
    \resizebox{0.8\textwidth}{!}{ %
    \begin{tabular}{llccccl}
    \toprule
    \textbf{Domain} & \textbf{Dataset} & \textbf{Train Size} & \textbf{Test Size} & \textbf{OOD} & \textbf{Evaluation Metric} \\
    \midrule
    \multirow{3}{*}{\textbf{Math}} 
        & \textsc{GSM8k}          & 1,500 & 1,319 & \xmark & \multirow{3}{*}{Numeric Match} \\
        & \textsc{Hendrycks-Math} & 1,500 & 500   & \xmark &  \\
        & \textsc{AIME-24/25}        & 0     & 60    & \cmark &  \\
    \midrule
    \multirow{2}{*}{\textbf{Code}} 
        & \textsc{LeetCode-Dataset}       & 2,585 & 223   & \xmark & \multirow{2}{*}{Unit Test Execution} \\
        & \textsc{MBPP}           & 0     & 974   & \cmark &  \\
    \midrule
    \multirow{2}{*}{\textbf{Translation}} 
        & \textsc{Flores-200}     & 3,000 & 300   & \xmark & \multirow{2}{*}{Hybrid (LLM-Judge + BLEU)} \\
        & \textsc{WMT-24}         & 0     & 500   & \cmark &  \\
    \midrule
    \textbf{General} 
        & \textsc{UltraChat}      & 2,500 & 200   & \xmark & Model-based (LLM-as-a-Judge) \\
    \bottomrule
    \end{tabular}
    }
    \vspace{-4mm}
\end{table}




% This formulation reduces sensitivity to reward model estimation errors by leveraging relative comparisons within output groups.

